\scanpage{37}
ouvertes de la partition \unsure{donnée|cherchée} de $Y$. Les feuilles de $X$ au dessus de celles-ci sont simplement les feuilles de la partition \inserted{privilégiée} canonique (intrinsèque) de $V=X|U$ (loc. sur $U$, c'est le produit de $U$ avec une feuille canonique d'une fibre). Si $\overline F_i$ sont ces feuilles, il faudra bien sûr que les $\Phi'_i=\overline F_i-V\cap\overline F_i=\overline F_i|Y'$ soient des réunions de feuilles de $X$.
\figure{\begin{tikzpicture}[x=1cm,y=1cm]\draw (.1,.65)--(0,.1) .. controls (-.1,-.08) and (.18,-.05) .. (.1,.1) -- (0,.16);\end{tikzpicture}}
\margin{\figure{\begin{tikzpicture}[x=1cm,y=1cm]\draw (.1,.65)--(0,.1) .. controls (-.1,-.08) and (.18,-.05) .. (.1,.1) -- (0,.16);\end{tikzpicture}} Il faut donc prouver que les $\overline F_i$ (donc les $\overline F'_i$) sont encore des ``espaces'' de notre Catégorie.}
Mais si $X'=X|Y'=X-V$ on est ramené à trouver la \deleted{\unsure{décomposition}}\inserted{partition canonique} pour $X'\unsure{-|\to}Y'$, mais $X'$ étant muni de plus des sous-espaces $\Phi'_i$ et de leurs intersections finies, tant qu'à faire. On gagne par hyp. de récurrence, à cela près qu'il aurait fallu partir non avec $X\unsure{-|\to}Y$, mais avec en plus une famille finie de sous-espaces fermés donnés de $X$. Pour paraphraser dans ce cas plus général la construction précédente, il faut simplement (remplaçant $X$ par $V$) pour un espace $X$ muni d'une famille \inserted{finie} de sous-espaces fermés, définir la \deleted{\unsure{sa}} partition privilégiée canonique (i.e. la plus grossière) pour ce système. \unsure{Or} la construction est probablement à faire en paraphrasant celle proposée dans la \inserted{$1^{\mathrm{re}}$} lettre à A Campo.
\mspage{38}
